\subsection{Árbol de recursión: divisiones de distinto tamaño}
\label{sec:act3-1}

\begin{enunciado}
Considere el siguiente algoritmo:
\begin{lstlisting}[language={}, numbers=none, frame=none, backgroundcolor={}]
Procedure RecursiveFilter(A, n)
    if n <= 1
        return

    for i = 1 to n
        Process(A[i])          // costo Theta(1)

    RecursiveFilter(A[1 ... floor(n/2)], floor(n/2))
    RecursiveFilter(A[1 ... floor(n/4)], floor(n/4))
\end{lstlisting}
\begin{enumerate}[label=(\alph*)]
  \item Formule la recurrencia de tiempo.
  \item Dibuje al menos los primeros tres niveles del árbol de recursión, indicando el
        tamaño de cada subproblema y el trabajo no recursivo de cada nodo.
  \item Determine una cota superior para el trabajo agregado de cada nivel.
  \item Utilizando el árbol, obtenga una cota $\Theta$ para el tiempo total.
  \item Explique por qué la forma estándar del Teorema Maestro no se aplica directamente a
        esta recurrencia.
\end{enumerate}
\end{enunciado}

% ============================================================
\subsubsection*{(a) Recurrencia de tiempo}

Llamamos $T(n)$ al tiempo que tarda \texttt{RecursiveFilter} con un arreglo de tamaño
$n$. Recorremos el pseudocódigo línea por línea.

Si $n \le 1$, el algoritmo solo hace la comparación \texttt{n <= 1} y termina. Eso cuesta
una constante, que llamamos $c$.

Si $n > 1$, el algoritmo primero hace la misma comparación, que cuesta una constante.
Después viene el for, que da $n$ vueltas. En cada vuelta se ejecuta \texttt{Process}, de
costo $\T{1}$, y se avanza el índice, así que el for cuesta $\T{n}$. Al final hay dos
llamadas recursivas: una con un arreglo de tamaño $\floor{n/2}$, que tarda
$T(\floor{n/2})$, y otra con un arreglo de tamaño $\floor{n/4}$, que tarda
$T(\floor{n/4})$. Suponemos que pasar el subarreglo $A[1 \dots k]$ no copia los datos,
sino que solo indica dónde termina, así que cuesta $\T{1}$.
La comparación y el for juntos cuestan a lo más $c\,n$ para alguna constante $c$, que
podemos tomar igual a la del caso base. Entonces:
\[
  T(n) =
  \begin{cases}
    c, & n \le 1, \\[4pt]
    T\!\left(\floor*{\dfrac{n}{2}}\right) + T\!\left(\floor*{\dfrac{n}{4}}\right) + c\,n, & n > 1.
  \end{cases}
\]
El término $T(\floor{n/2})$ viene de la primera llamada recursiva, $T(\floor{n/4})$ de la
segunda, y $c\,n$ del for y la comparación. En adelante omitimos los pisos y escribimos
\[
  T(n) = T\!\left(\frac{n}{2}\right) + T\!\left(\frac{n}{4}\right) + c\,n.
\]
Así los tamaños del árbol quedan como $n/2$, $n/4$, $n/8$, y la suma de cada nivel es
fácil de calcular. Esto no cambia el resultado. El piso solo achica los tamaños, así que
la cota superior que obtenemos sin pisos también vale con ellos, y la cota inferior solo
usa la raíz, que tiene tamaño $n$.

% ============================================================
\subsubsection*{(b) Árbol de recursión}

Cada nodo del árbol es una llamada. Dentro de cada nodo anotamos el tamaño del
subproblema y, debajo, su trabajo no recursivo, que es $c$ por el tamaño. Cada nodo de
tamaño $s$ tiene dos hijos, de tamaños $s/2$ y $s/4$. En la tabla de abajo va la suma de cada
nivel.

\begin{center}
\begin{forest}
  for tree={
    draw, rounded corners, align=center, font=\small,
    l sep=12mm, s sep=4mm, inner sep=3pt, edge={->}
  }
  [{$n$\\[1pt]$c\,n$}
    [{$n/2$\\[1pt]$c\,n/2$}
      [{$n/4$\\[1pt]$c\,n/4$}]
      [{$n/8$\\[1pt]$c\,n/8$}]
    ]
    [{$n/4$\\[1pt]$c\,n/4$}
      [{$n/8$\\[1pt]$c\,n/8$}]
      [{$n/16$\\[1pt]$c\,n/16$}]
    ]
  ]
\end{forest}
\end{center}

\begin{center}
\begin{tabular}{ccl}
  \toprule
  nivel & tamaños de los nodos & trabajo del nivel \\
  \midrule
  0 & $n$ & $c\,n$ \\[2pt]
  1 & $\frac{n}{2},\ \frac{n}{4}$ & $c\left(\frac12 + \frac14\right) n = \frac34\, c\,n$ \\[4pt]
  2 & $\frac{n}{4},\ \frac{n}{8},\ \frac{n}{8},\ \frac{n}{16}$
    & $c\left(\frac14 + \frac18 + \frac18 + \frac1{16}\right) n = \frac{9}{16}\, c\,n
       = \left(\frac34\right)^2 c\,n$ \\
  \bottomrule
\end{tabular}
\end{center}

La rama de más a la izquierda siempre divide por $2$, así que llega a tamaño $1$ después
de $\log_2 n$ niveles. La de más a la derecha siempre divide por $4$, así que llega antes,
después de $\log_4 n$ niveles. Por eso el árbol no es completo: a partir del nivel
$\log_4 n$ algunas ramas ya terminaron.

% ============================================================
\subsubsection*{(c) Cota superior para el trabajo de cada nivel}

En la tabla se ve que el trabajo de cada nivel es $\frac34$ del anterior. La razón es que
un nodo de tamaño $s$ tiene hijos de tamaños $s/2$ y $s/4$, que suman
\[
  \frac{s}{2} + \frac{s}{4} = \frac34\, s.
\]
Es decir, al bajar un nivel, la suma de los tamaños se multiplica por $\frac34$. Como el
trabajo de cada nodo es $c$ por su tamaño, el trabajo del nivel también se multiplica
por $\frac34$. Cuando una rama termina, sus nodos dejan de aparecer en los niveles
siguientes, y el redondeo hacia abajo solo achica los tamaños, así que el trabajo de un
nivel nunca supera ese valor. Por lo tanto, el trabajo del nivel $i$ cumple
\[
  \text{trabajo del nivel } i \le \left(\frac34\right)^{i} c\,n.
\]

% ============================================================
\subsubsection*{(d) Cota $\Theta$ para el tiempo total}

\textbf{Cota superior.} El tiempo total es la suma del trabajo de todos los niveles.
Usando la cota de (c) y dejando que la suma siga para siempre:
\[
  T(n) \le c\,n + \frac34\, c\,n + \left(\frac34\right)^2 c\,n + \cdots
       = c\,n \left( 1 + \frac34 + \left(\frac34\right)^2 + \cdots \right).
\]
Lo que queda entre paréntesis es una serie geométrica de razón $3/4$:
\[
  1 + \frac34 + \left(\frac34\right)^2 + \cdots = \frac{1}{1 - 3/4} = 4.
\]
Reemplazando este valor, el trabajo de los nodos internos es a lo más $4c\,n$.

Falta sumar las hojas, que son las llamadas con $n \le 1$. Se suman aparte porque una
hoja no entra al for. Solo hace la comparación \texttt{n <= 1}, así que cuesta $c$ aunque
su tamaño sea $0$, y la cota de (c), que cuenta $c$ por el tamaño, no la incluye bien.

Cada nodo interno tiene tamaño al menos $2$, así que su trabajo es al menos $2c$. Como el
trabajo de todos los nodos internos suma a lo más $4c\,n$, hay a lo más $2n$ nodos
internos. Cada nodo interno tiene dos hijos, así que hay a lo más $4n$ hojas, y entre
todas cuestan a lo más $4c\,n$. Por lo tanto,
\[
  T(n) \le 4c\,n + 4c\,n = 8c\,n.
\]

\textbf{Cota inferior.} La raíz sola ya hace un trabajo de $c\,n$, y los demás nodos no
pueden restar. Por lo tanto, $T(n) \ge c\,n$.

Juntando ambas cotas:
\[
  c \cdot n \le T(n) \le 8c \cdot n.
\]
Según la definición vista en clases, $T(n) \in \T{n}$ si existen dos constantes
$c_1, c_2 > 0$ y un $n_0 \in \N$ tales que $0 \le c_1 n \le T(n) \le c_2 n$ para todo
$n \ge n_0$. Luego, $c_1 = c$ y $c_2 = 8c$, y las dos cotas valen desde $n = 2$, así que
$n_0 = 2$. Por lo tanto, $T(n) \in \T{n}$.

El trabajo total es del mismo orden que el de la raíz. Como cada nivel hace $\frac34$ del
trabajo del anterior, el trabajo decrece geométricamente y la raíz domina, igual que en
el ejemplo de clases con $T(n) = 3T(n/4) + c\,n^2$.

% ============================================================
\subsubsection*{(e) Por qué no se aplica el Teorema Maestro}

La forma estándar del Teorema Maestro es para recurrencias del tipo
\[
  T(n) = a\, T\!\left(\frac{n}{b}\right) + f(n),
\]
donde los $a$ subproblemas tienen todos el mismo tamaño $n/b$. En nuestra recurrencia hay
dos subproblemas de tamaños distintos, $n/2$ y $n/4$, así que no existe un único $b$ que
la describa. Por eso no se puede comparar $f(n)$ con $n^{\log_b a}$, que es lo que hace
el teorema. Este es justamente uno de los casos que se vieron en clases como
``subproblemas de tamaños distintos'', junto con el ejemplo $T(n/3) + T(2n/3) + n$, que
también se resuelve con el árbol de recursión.

